
\chapter{2021年天津卷}

\section{单选题}

\begin{problem}
设集合 \( A=\{-1,0,1\} \),\( B=\{1,3,5\} \),\( C=\{0,2,4\} \),则 \( \left( A\cap B \right)\cup C= \)(\quad)

\choices
    {\(\{0\}\)}
    {\(\{0,1,3,5\}\)}
    {\(\{0,1,2,4\}\)}
    {\(\{0,2,3,4\}\)}
\end{problem}

\begin{problem}
已知 \( a\in \R \),则``\( a>6 \)''是``\( a^2>36 \)''的(\quad)

\choices
    {充分不必要条件}
    {必要不充分条件}
    {充要条件}
    {既不充分也不必要条件}
\end{problem}

\begin{problem}
函数 \( y=\dfrac{\ln |x|}{x^2+2} \) 的图像大致为(\quad)

\choices
    {\choicebitmap[width=0.15\paperwidth]{img/2021/tianjin/q3_optA.png}}
    {\choicebitmap[width=0.15\paperwidth]{img/2021/tianjin/q3_optB.png}}
    {\choicebitmap[width=0.15\paperwidth]{img/2021/tianjin/q3_optC.png}}
    {\choicebitmap[width=0.15\paperwidth]{img/2021/tianjin/q3_optD.png}}
\end{problem}

\begin{problem}
从某网络平台推荐的影视作品中抽取 \( 400 \) 部,统计其评分分布数据,将所得 \( 400 \) 个评分数据分为 \( 8 \) 组:
\( [66,70) \),\( [70,74) \),\(\cdots\),\( [94,98] \),并整理得到如下的频率分布直方图,则评分在区间 \( [82,86) \) 内的影视作品数量是(\quad)

\bitmapfigure[max width=5.5cm]{img/2021/tianjin/q4_fig1.png}

\choices
    {20}
    {40}
    {64}
    {80}
\end{problem}

\begin{problem}
设 \( a=\log_{2}0.3 \),\( b=\log_{\frac{1}{2}}0.4 \),\( c=0.4^{0.3} \),则 \(a\),\(b\),\(c\) 的大小关系为(\quad)

\choices
    {\(a<b<c\)}
    {\(c<a<b\)}
    {\(b<c<a\)}
    {\(a<c<b\)}
\end{problem}

\begin{problem}
两个圆锥的底面是一个球的同一截面,顶点均在球面上,若球的体积为 \( \dfrac{32\pi}{3} \),两个圆锥的高之比为 \( 1:3 \),则这两个圆锥的体积之和为(\quad)

\choices
    {\(3\pi\)}
    {\(4\pi\)}
    {\(9\pi\)}
    {\(12\pi\)}
\end{problem}

\begin{problem}
若 \( 2^a=5^b=10 \),则 \( \dfrac{1}{a}+\dfrac{1}{b}= \)(\quad)

\choices
    {\(-1\)}
    {\(\lg 7\)}
    {\(1\)}
    {\(\log_{7}10\)}
\end{problem}

\begin{problem}
已知双曲线 \( \dfrac{x^2}{a^2}-\dfrac{y^2}{b^2}=1 \)\( \left( a>0,b>0 \right) \) 的右焦点与抛物线 \( y^2=2px \)\( \left( p>0 \right) \) 的焦点重合,抛物线的准线交双曲线于 \(A\),\(B\) 两点,交双曲线的渐近线于 \(C\),\(D\) 两点,若 \( |CD|=\sqrt2\,|AB| \),则双曲线的离心率为(\quad)

\choices
    {\(\sqrt2\)}
    {\(\sqrt3\)}
    {\(2\)}
    {\(3\)}
\end{problem}

\begin{problem}
设 \( a\in \R \),函数
\( f(x)=\begin{cases}
\cos\left( 2\pi x-2\pi a \right), & x<a,\\
x^2-2(a+1)x+a^2+5, & x\ge a
\end{cases} \)
若 \( f(x) \) 在区间 \( \left( 0,+\infty \right) \) 内恰有 \( 6 \) 个零点,则 \( a \) 的取值范围是(\quad)

\choices
    {\(\left( 2,\frac{9}{4} \right]\cup \left( \frac{5}{2},\frac{11}{4} \right]\)}
    {\(\left( \frac{7}{4},2 \right)\cup \left( \frac{5}{2},\frac{11}{4} \right]\)}
    {\(\left( 2,\frac{9}{4} \right]\cup \left[ \frac{11}{4},3 \right)\)}
    {\(\left( \frac{7}{4},2 \right)\cup \left[ \frac{11}{4},3 \right)\)}
\end{problem}

\section{填空题}

\begin{problem}
\( \i \) 是虚数单位,复数 \( \dfrac{9+2\i}{2+\i}=\fillinblank{} \).
\end{problem}

\begin{problem}
在 \( \left( 2x^3+\dfrac{1}{x} \right)^6 \) 的展开式中,\( x^6 \) 的系数是 \(\fillinblank{}\).
\end{problem}

\begin{problem}
若斜率为 \( \sqrt3 \) 的直线与 \( y \) 轴交于点 \(A\),与圆 \( x^2+\left( y-1 \right)^2=1 \) 相切于点 \(B\),则 \( |AB|=\fillinblank{} \).
\end{problem}

\begin{problem}
若 \( a>0 \),\( b>0 \),则 \( \dfrac{1}{a}+\dfrac{a}{b^2}+b \) 的最小值为 \(\fillinblank{}\).
\end{problem}

\begin{problem}
甲,乙两人在每次猜谜活动中各猜一个谜语,若一方猜对且另一方猜错,则猜对的一方获胜,否则本次平局. 已知每次活动中,甲,乙猜对的概率分别为 \( \dfrac{5}{6} \),\( \dfrac{3}{5} \),且每次活动中甲,乙猜对与否互不影响,各次活动也互不影响,则一次活动中,甲获胜的概率为 \(\fillinblank{}\),\( 3 \) 次活动中,甲至少获胜 \( 2 \) 次的概率为 \(\fillinblank{}\).
\end{problem}





\begin{problem}
在边长为 \( 1 \) 的等边三角形 \(ABC\) 中,\(D\) 为线段 \(BC\) 上的动点,\( DE\perp AB \) 且交 \(AB\) 于点 \(E\),\( DF\parallel AB \) 且交 \(AC\) 于点 \(F\),则 \( \left| 2\overrightarrow{BE}+\overrightarrow{DF} \right| \) 的值为 \(\fillinblank{}\);\( \left( \overrightarrow{DE}+\overrightarrow{DF} \right)\cdot \overrightarrow{DA} \) 的最小值为 \(\fillinblank{}\).

\bitmapfigure[max width=6.0cm]{img/2021/tianjin/q15_fig1.png}
\end{problem}




\section{解答题}

\begin{problem}
在 \( \triangle ABC \) 中,角 \(A\),\(B\),\(C\) 所对的边分别为 \( a \),\( b \),\( c \),已知 \( \sin A:\sin B:\sin C=2:1:\sqrt2 \),\( b=\sqrt2 \).

\begin{enumerate}
\item 求 \( a \) 的值;
\item 求 \( \cos C \) 的值;
\item 求 \( \sin\left( 2C-\dfrac{\pi}{6} \right) \) 的值.
\end{enumerate}
\end{problem}

\begin{problem}
如图,在棱长为 \( 2 \) 的正方体 \( ABCD-A_1B_1C_1D_1 \) 中,\(E\) 为棱 \(BC\) 的中点,\(F\) 为棱 \(CD\) 的中点.

\begin{enumerate}
\item 求证:\( D_1F\parallel \) 平面 \(A_1EC_1\);
\item 求直线 \(AC_1\) 与平面 \(A_1EC_1\) 所成角的正弦值;
\item 求二面角 \( A-A_1C_1-E \) 的正弦值.
\end{enumerate}

\bitmapfigure[max width=5cm]{img/2021/tianjin/q17_fig1.png}
\end{problem}

\begin{problem}
已知椭圆 \( \dfrac{x^2}{a^2}+\dfrac{y^2}{b^2}=1 \)\( \left( a>b>0 \right) \) 的右焦点为 \(F\),上顶点为 \(B\),离心率为 \( \dfrac{2\sqrt5}{5} \),且 \( |BF|=\sqrt5 \).

\begin{enumerate}
\item 求椭圆的方程;
\item 直线 \( l \) 与椭圆有唯一的公共点 \(M\),与 \( y \) 轴的正半轴交于点 \(N\),过 \(N\) 与 \(BF\) 垂直的直线交 \( x \) 轴于点 \(P\). 若 \( MP\parallel BF \),求直线 \( l \) 的方程.
\end{enumerate}
\end{problem}

\begin{problem}
已知 \( \{a_n\} \) 是公差为 \( 2 \) 的等差数列,其前 \( 8 \) 项和为 \( 64 \). \( \{b_n\} \) 是公比大于 \( 0 \) 的等比数列,\( b_1=4 \),\( b_3-b_2=48 \).

\begin{enumerate}
\item 求 \( \{a_n\} \) 和 \( \{b_n\} \) 的通项公式;
\item 记 \( c_n=b_{2n}+\dfrac{1}{b_n} \),\( n\in \N^* \),
\begin{enumerate}
\item 证明 \( \{c_n^2-c_{2n}\} \) 是等比数列;
\item 证明 \( \sum_{k=1}^n \sqrt{\dfrac{a_ka_{k+1}}{c_k^2-c_{2k}}}<2\sqrt2 \)\( \left( n\in \N^* \right) \).
\end{enumerate}
\end{enumerate}
\end{problem}

\begin{problem}
已知 \( a>0 \),函数 \( f(x)=ax-x\e^x \).

\begin{enumerate}
\item 求曲线 \( y=f(x) \) 在点 \( \left( 0,f(0) \right) \) 处的切线方程;
\item 证明 \( f(x) \) 存在唯一的极值点;
\item 若存在 \( a \),使得 \( f(x)\le a+b \) 对任意 \( x\in \R \) 成立,求实数 \( b \) 的取值范围.
\end{enumerate}
\end{problem}

